\documentclass[10pt,a4paper]{amsart}
\usepackage[margin=19mm]{geometry}
\usepackage{amsmath,amssymb,mathtools}
\usepackage[scr=rsfs,cal=euler]{mathalfa}
\usepackage{xfrac}
\usepackage[unicode,hidelinks,bookmarks=false]{hyperref}
\newcommand{\ie}{i.e.\ }
\newcommand{\cf}{cf.\ }
\newcommand{\resp}{respectively\ }
\newcommand{\Subsection}{\subsection}
\providecommand{\nobreakdash}{\nobreak\hbox{-}}
\setlength{\emergencystretch}{2em}
\multlinegap0pt
\usepackage{bm}
\usepackage{bbm}
\usepackage{dsfont}
\usepackage{xspace}
\usepackage{cancel}
\usepackage{mathtools}
\usepackage{amsthm}
\makeatletter
\@ifundefined{theorem}{\newtheorem{theorem}{Theorem}}{}
\@ifundefined{claim}{\newtheorem{claim}[theorem]{Claim}}{}
\@ifundefined{fact}{\newtheorem{fact}[theorem]{Fact}}{}
\makeatother
\newcommand{\forkindep}[1][]{%
  \mathrel{
    \mathop{
      \vcenter{
        \hbox{\oalign{\noalign{\kern-.3ex}\hfil$\vert$\hfil\cr
              \noalign{\kern-.7ex}
              $\smile$\cr\noalign{\kern-.3ex}}}
      }
    }\displaylimits_{#1}
  }
}
\newcommand{\gs}{\sigma}
\newcommand{\gw}{\omega}
\newcommand{\vd}{\mathrm{OD}}
\title[More on the Boolean Prime Ideal Theorem]{More on the Boolean Prime Ideal Theorem}
\author{Jacob Kowalczyk, Jindrich Zapletal}
\date{Source: arXiv:2607.16747v1 (CC BY 4.0)}
\begin{document}
\maketitle

\noindent\emph{Source and licence.} More on the Boolean Prime Ideal Theorem. Version arXiv:2607.16747v1 (CC BY 4.0), distributed under the Creative Commons Attribution 4.0 International licence, as recorded in the arXiv licence metadata for this exact version and shown on the abstract page https://arxiv.org/abs/2607.16747v1. Full rights notice: licenses/arxiv-2607.16747-CC-BY-4.0.txt.\par\smallskip

\noindent\textit{Editorial scope.} Counted unit: the Theorem whose source label is \texttt{ksigmacriterion} in the article (source0.tex lines 412--421, source label \texttt{ksigmacriterion}), delivered together with the article's own complete original proof of it (source0.tex lines 423--448, one \texttt{proof} environment of 26 lines). No mathematics is written, repaired or re-derived: statement and proof are the article's own text line for line. Required context retained uncounted: facticfact (lines 747--753). Every definition, display and label of those lines is retained unchanged. Omitted from this package and not redistributed: 1--411, 422--422, 449--746, 754--811. Nothing inside a retained excerpt is omitted or altered: every retained body is delivered exactly as the article's lines are written, so no omission receipt of any kind accompanies this package. Citations: the retained excerpts carry no citation command at all, so no bibliography entry of the article is transcribed and no reference is redistributed. Reference adaptation: none was needed and none was made; every reference inside the delivered unit resolves inside it, so no unresolved reference or citation is produced. Printed slips kept literal: no printed word, symbol, hypothesis or number of the counted statement or of its proof was altered. Layout and class: the article's own class is \texttt{article} with packages including amsmath, amsthm, amssymb, amscd; the wrapper is the lane's pinned \texttt{amsart} editorial wrapper at 10pt on A4, which restates the theorem-like environments the retained text needs and adds the article's own macro definitions that the retained text needs (\texttt{\textbackslash forkindep}, \texttt{\textbackslash gs}, \texttt{\textbackslash gw}, \texttt{\textbackslash vd}), with extra wrapper packages (bm, bbm, dsfont, xspace, cancel, mathtools). The delivered numbering is the unit's own. Assets: no retained span contains a figure environment, a graphics-inclusion command, a PDF-page inclusion, a table or a TikZ picture; no image file is copied into this package, the package declares no asset dependency and no image file is redistributed with it. Licence: version arXiv:2607.16747v1 of this article is distributed under the Creative Commons Attribution 4.0 International licence, as recorded in the arXiv licence metadata for this exact version and shown on the abstract page https://arxiv.org/abs/2607.16747v1. A full rights notice is delivered at licenses/arxiv-2607.16747-CC-BY-4.0.txt. Characters: every retained excerpt is pure ASCII, so the delivered engine, XeLaTeX with its default Latin Modern text font, reports no missing character.
	\begin{fact}
		\label{facticfact}
		\begin{enumerate}
			\item Let $A$ be a set and $z$ a set of ordinals such that $A\in\vd_z$. Then $A$ is well-orderable iff $A\subset\vd_z$.
			\item if $y_0\forkindep_z y_1$ then $\vd_{zy_0}\cap\vd_{zy_1}=\vd_z$.
			\item if $y_0\forkindep_z y_1$ and $A_0\in \vd_{zy_0}$ and $A_1\in\vd_{zy_1}$ are disjoint subsets of $\vd_z$ then there is $B\in\vd_z$ such that $A_0\subseteq B$ and $B\cap A_1=0$.	\end{enumerate}
	\end{fact}

\begin{theorem}
\label{ksigmacriterion}
Let $T$ be a deductively closed $K_\gs$-theory on a locally compact space $X$. Suppose that

\begin{enumerate}
\item[(*)] there are compact sets $U_n$ for $n\in\gw$ such that  $T=\bigcup_nU_n$ and for every $n\in\gw$ and all compact sets $K_0, K_1\subset X^*$, if for all $\phi_0\in K_0$ and $\phi_1\in K_1$ the implication $\phi_0\to\phi_1$ belongs to $U_n$, then there is a formula $\psi$ such that for all $\phi_0\in K_0$ and $\phi_1\in K_1$ the implications $\phi_0\to\psi$ and $\psi\to\phi_1$ belong to $T$.
\end{enumerate}

\noindent Then $T$ is balanced.
\end{theorem}

\begin{proof}
Suppose that $z, y_0, y_1$ are sets of ordinals such that $X, T, \langle U_n\colon n\in\gw\rangle\in\vd_z$ and $y_0\forkindep_zy_1$, and suppose that $\phi_0\in\vd_{zy_0}$ and $\phi_1\in\vd_{zy_1}$ are formulas in $X^*$ such that the implication $\phi_0\to\phi_1$ belongs to $T$. There must be a number $n$ such that the implication belongs to $U_n$.





\begin{claim}
There is a closed set $K_0\subseteq X^*$ in $\vd_z$ such that $\phi_0\in K_0$ and for every $\psi_0\in K_0$, the implication $\psi_0\to\phi_1$ belongs to $U_n$.
\end{claim}

\begin{proof}
Let $W_0$ be the set of all basic open subsets of $X^*$ which contain $\phi_0$, and let $W_1$ be the set of all basic open subsets of $X^*$ which contain no formula $\phi'_0$ such that $\phi'_0\to\phi_1\in U_n$. These are obviously disjoint subsets of $\vd_z$, $W_0\in\vd_{zy_0}$ and $W_1\in\vd_{zy_1}$, so by Fact~\ref{facticfact}(3) there must be a set $B\in\vd_z$ consisting of basic open subsets of $X^*$ such that $A_0\cap B=0$ and $A_1\subseteq B$. Write $K_0=X\setminus \bigcup B\in\vd_z$. As $B$ is disjoint from $A_0$, it is clear that $\phi_0\in K_0$. As $A_1\subseteq B$ and the set $U_n$ is closed, it follows that for every formula $\psi_0\in K_0$, $\psi_0\to\phi_1\in U_n$ as desired.
\end{proof}

\noindent A symmetric argument yields

\begin{claim}
There is a closed set $K_1\subseteq X^*$ in $\vd_z$ such that $\phi_1\in K_1$ and for every $\psi_1\in K_1$, the implication $\phi_0\to\psi_1$ belongs to $U_n$.
\end{claim}


\noindent Use the locally compact assumption to shrink the sets $K_0, K_1$ to compact. Further shrinking $K_0$ to the compact set $\{\psi_0\in K_0\colon\forall \psi_1\in K_1\ \psi_0\to\psi_1\in U_n\}$  and $K_1$ to the compact set $\{\psi_1\in K_1\colon\forall \psi_0\in K_0\ \psi_0\to\psi_1\in U_n\}$, we may achieve the situation where all implications from an element of $K_0$ to an element of $K_1$ belong to the set $U_n$.

Now, use the assumption (*) to conclude that there is a formula $\psi$ such that for all $\psi_0\in K_0$ and all $\psi_1\in K_1$, the implications $\psi_0\to\psi$ and $\psi\to\psi_1$ belong to $T$. This is a ${\mathbf\Sigma}^1_2$-statement; by Shoenfield's absoluteness, such a formula $\psi$ must exist in $\vd_z$. The balance of the theory $T$ follows.
\end{proof}

\end{document}
